\documentclass[10pt,a4paper]{amsart}
\usepackage[margin=19mm]{geometry}
\usepackage{amsmath,amssymb,mathtools}
\usepackage[scr=rsfs,cal=euler]{mathalfa}
\usepackage{xfrac}
\usepackage[unicode,hidelinks,bookmarks=false]{hyperref}
\newcommand{\ie}{i.e.\ }
\newcommand{\cf}{cf.\ }
\newcommand{\resp}{respectively\ }
\newcommand{\Subsection}{\subsection}
\providecommand{\nobreakdash}{\nobreak\hbox{-}}
\setlength{\emergencystretch}{2em}
\multlinegap0pt
\usepackage{amssymb}
\newtheorem{theorem}{Theorem}
\newcommand{\dis}{\displaystyle}
\title{Picture Fuzzy Multigroups}
\author{Taiwo O. Sangodapo}
\date{Source: arXiv:2607.06580v1 (CC BY 4.0)}
\begin{document}
\maketitle

\noindent\emph{Source and licence.} Picture Fuzzy Multigroups. Version arXiv:2607.06580v1 (CC BY 4.0), distributed by arXiv under the Creative Commons Attribution 4.0 International licence, http://creativecommons.org/licenses/by/4.0/, as recorded in the arXiv licence metadata for this exact version and shown on the abstract page https://arxiv.org/abs/2607.06580v1. Full licence notice: \texttt{licenses/arxiv-2607.06580-CC-BY-4.0.txt}.\par\smallskip

\noindent\textit{Editorial scope.} Counted unit: the article's theorem theorem (source0.tex lines 175--199). The article's complete original proof of it is delivered with it (source0.tex lines 201--306, one \texttt{proof} environment of 106 lines). No mathematics is written, repaired or re-derived: the statement and the proof are the article's own lines, in the article's own notation, and no retained line is substituted or re-flowed.  No further source-local context is needed: every cross-reference of every retained line resolves inside the two delivered spans.  Nothing was omitted from any retained span, so the package carries no omission record.  No retained line contains a citation, so the wrapper carries no bibliography and no bibliography entry.  No retained line contains a figure environment, an image-inclusion command or drawing code, so no image is delivered and the package declares no asset dependency.  Editorial formatting only, in addition: an AMS \texttt{amsart} preamble that restates the article's own declarations and macros used by the retained lines, namely the environments \texttt{theorem} on separate counters, together with the standard packages those lines need.  The retained environments print in this document as Theorem 1 in order of appearance, whereas the article prints the same items with the numbering of its own full document.  The complete native source of the article as distributed in the arXiv e-print for this exact version is bundled unchanged as \texttt{sources/204652/source0.tex}.
\begin{theorem}
Let $D, E, D_i \in PFMS(C).$ Then, the following properties hold:
\begin{itemize}
\item [i.] $(D^{-1})^{-1} = D$
\item [ii.] $D \subseteq E~\Rightarrow~D^{-1} \subseteq E^{-1}$
\item [iii.] $ \dis \left(\bigcup^n_{i = 1} D_i\right)^{-1} = \bigcup^n_{i = 1} D_i^{-1}$
\item [iv.] $\dis \left(\bigcap^n_{i = 1} D_i\right)^{-1} = \bigcap^n_{i = 1} D_i^{-1}$
\item [v.] $(D \circ E) = E^{-1} \circ D^{-1}$
\item [vi.] 
\begin{eqnarray*}
PCM_{D \circ E}(g_1) & = & \bigvee_{g_2 \in C} \lbrace PCM_D(g_2) \wedge PCM_E(g_2^{-1}g_1)\rbrace~~ \forall~~ g_1 \in C\\
& = & \bigvee_{g_2 \in C} \lbrace PCM_D(g_1 g^{-1}_2) \wedge PCM_E(g_2)\rbrace~~ \forall~~ g_1 \in C
\end{eqnarray*}

\begin{eqnarray*}
NeCM_{D \circ E}(g_1) & = & \bigvee_{g_2 \in C} \lbrace NeCM_D(g_2) \wedge NeCM_E(g_2^{-1}g_1)\rbrace~~ \forall~~ g_1 \in C\\
& = & \bigvee_{g_2 \in C} \lbrace NeCM_D(g_1 g^{-1}_2) \wedge NeCM_E(g_2)\rbrace~~ \forall~~ g_1 \in C
\end{eqnarray*}
and 
\begin{eqnarray*}
NCM_{D \circ E}(g_1) & = & \bigwedge_{g_2 \in C} \lbrace NCM_D(g_2) \vee NeCM_E(g_2^{-1}g_1)\rbrace~~ \forall~~ g_1 \in C\\
& = & \bigwedge_{g_2 \in C} \lbrace NCM_D(g_1 g^{-1}_2) \vee NCM_E(g_2)\rbrace~~ \forall~~ g_1 \in C
\end{eqnarray*}
\end{itemize} 
\end{theorem}

\begin{proof}
i.
\begin{eqnarray*}
PCM_{(D^{-1})^{-1}}(g_1) & = & PCM_{(D^{-1})}(g^{-1})\\
& = & PCM_D(g^{-1})^{-1} \\
& = & PCM_D(g)~~ \forall~~ g \in C.~~ \text{Since}~~ C ~ \text{is~~ a~~ group}~~((g)^{-1})^{-1} = g
\end{eqnarray*}

\begin{eqnarray*}
NeCM_{(D^{-1})^{-1}}(g_1) & = & NeCM_{(D^{-1})}(g^{-1})\\
& = & NeCM_D(g^{-1})^{-1} \\
& = & NeCM_D(g)~~ \forall~~ g \in C.~~ \text{Since}~~ C ~ \text{is~~ a~~ group}~~((g)^{-1})^{-1} = g
\end{eqnarray*} and 

\begin{eqnarray*}
NCM_{(D^{-1})^{-1}}(g_1) & = & NCM_{(D^{-1})}(g^{-1})\\
& = & NCM_D(g^{-1})^{-1} \\
& = & NCM_D(g)~~ \forall~~ g \in C.~~ \text{Since}~~ C ~ \text{is~~ a~~ group}~~((g)^{-1})^{-1} = g
\end{eqnarray*}
Hence, $D = (D^{-1})^{-1}$\\

ii. Given that $D \subseteq E,$ this implies that, 
$$PCM_{D}(g^{-1}) \leq PCM_{E}(g^{-1})~~ \forall~~ g \in C$$
$$PCM_{D^{-1}}(g) \leq PCM_{E}(g),$$

$$NeCM_{D}(g^{-1}) \leq NeCM_{E}(g^{-1})~~ \forall~~ g \in C$$
$$NeCM_{D^{-1}}(g) \leq NeCM_{E}(g)$$ and 

$$NCM_{D}(g^{-1}) \geq NCM_{E}(g^{-1})~~ \forall~~ g \in C$$
$$NCM_{D^{-1}}(g) \geq NCM_{E}(g)$$ So, $D^{-1} \subseteq E^{-1}$\\

iii. \begin{eqnarray*}
PCM_{\dis \left(\bigcup^n_{i = 1} D_i\right)^{-1}}(g) & = & PCM_{\dis \left(\bigcup^n_{i = 1} D_i\right)}(g^{-1})\\
& = & \bigvee \lbrace  PCM_{D_i}(g^{-1});~~ i = 1, 2, \cdots, n \rbrace~~ \text{by~Definition ?} \\
& = & \bigvee \lbrace  PCM_{D_i}^{-1}(g);~~ i = 1, 2, \cdots, n \rbrace\\
& = & PCM_{\dis \left(\bigcup^n_{i = 1} {D_i}^{-1}\right)}(g)~~ \text{by~Definition ?}
\end{eqnarray*}

\begin{eqnarray*}
NeCM_{\dis \left(\bigcup^n_{i = 1} D_i\right)^{-1}}(g) & = & NeCM_{\dis \left(\bigcup^n_{i = 1} D_i\right)}(g^{-1})\\
& = & \bigvee \lbrace  NeCM_{D_i}(g^{-1});~~ i = 1, 2, \cdots, n \rbrace~~ \text{by~Definition ?} \\
& = & \bigvee \lbrace  NeCM_{D_i}^{-1}(g);~~ i = 1, 2, \cdots, n \rbrace\\
& = & NeCM_{\dis \left(\bigcup^n_{i = 1} {D_i}^{-1}\right)}(g)~~ \text{by~Definition ?}
\end{eqnarray*}
and

\begin{eqnarray*}
NCM_{\dis \left(\bigcap^n_{i = 1} D_i\right)^{-1}}(g) & = & NCM_{\dis \left(\bigcap^n_{i = 1} D_i\right)}(g^{-1})\\
& = & \bigwedge \lbrace  NCM_{D_i}(g^{-1});~~ i = 1, 2, \cdots, n \rbrace~~ \text{by~Definition ?} \\
& = & \bigwedge \lbrace  NCM_{D_i}^{-1}(g);~~ i = 1, 2, \cdots, n \rbrace\\
& = & NCM_{\dis \left(\bigcap^n_{i = 1} {D_i}^{-1}\right)}(g)~~ \text{by~Definition ?}
\end{eqnarray*}
Therefore, $ \dis \left(\bigcup^n_{i = 1} D_i\right)^{-1} = \bigcup^n_{i = 1} D_i^{-1}.$\\

iv. Property iv can be done in a similar manner.\\

v.
\begin{eqnarray*}
PCM_{ (D \circ E)^{-1}}(g_1) & = & PCM_{ (D \circ E)}(g^{-1}_1)\\
& = & \bigvee \lbrace  PCM_{D}(g_2) \wedge PCM_{E}(g_3);~~ g_2, g_3 \in C,~~ g_2g_3 = g_1^{-1} \rbrace \\
& = & \bigvee \lbrace  PCM_{E}(g_3) \wedge PCM_{D}(g_2);~~ g_2, g_3 \in C,~~ (g_2g_3)^{-1} = g_1 \rbrace\\
& = & \bigvee \lbrace  PCM_{E}(g^{-1}_3)^{-1} \wedge PCM_{D}(g^{-1}_2)^{-1};~~ g_2, g_3 \in C,~~ (g_2g_3)^{-1} = g_1 \rbrace\\
& = & \bigvee \lbrace  PCM_{E^{-1}}(g^{-1}_3) \wedge PCM_{D^{-1}}(g^{-1}_2);~~ g^{-1}_2, g^{-1}_3 \in C,~~ g^{-1}_2g^{-1}_3 = g_1 \rbrace\\
& = & PCM_{E^{-1} \circ D^{-1}}(g_1)~~ \forall~~ g_1 \in C.
\end{eqnarray*}

\begin{eqnarray*}
NeCM_{ (D \circ E)^{-1}}(g_1) & = & NeCM_{ (D \circ E)}(g^{-1}_1)\\
& = & \bigvee \lbrace  NeCM_{D}(g_2) \wedge NeCM_{E}(g_3);~~ g_2, g_3 \in C,~~ g_2g_3 = g_1^{-1} \rbrace \\
& = & \bigvee \lbrace  NeCM_{E}(g_3) \wedge NeCM_{D}(g_2);~~ g_2, g_3 \in C,~~ (g_2g_3)^{-1} = g_1 \rbrace\\
& = & \bigvee \lbrace  NeCM_{E}(g^{-1}_3)^{-1} \wedge NeCM_{D}(g^{-1}_2)^{-1};~~ g_2, g_3 \in C,~~ (g_2g_3)^{-1} = g_1 \rbrace\\
& = & \bigvee \lbrace  NeCM_{E^{-1}}(g^{-1}_3) \wedge NeCM_{D^{-1}}(g^{-1}_2);~~ g^{-1}_2, g^{-1}_3 \in C,~~ g^{-1}_2g^{-1}_3 = g_1 \rbrace\\
& = & NeCM_{E^{-1} \circ D^{-1}}(g_1)~~ \forall~~ g_1 \in C.
\end{eqnarray*}
and 

\begin{eqnarray*}
NCM_{ (D \circ E)^{-1}}(g_1) & = & NCM_{ (D \circ E)}(g^{-1}_1)\\
& = & \bigwedge \lbrace  NCM_{D}(g_2) \vee NCM_{E}(g_3);~~ g_2, g_3 \in C,~~ g_2g_3 = g_1^{-1} \rbrace \\
& = & \bigwedge \lbrace  NCM_{E}(g_3) \vee NCM_{D}(g_2);~~ g_2, g_3 \in C,~~ (g_2g_3)^{-1} = g_1 \rbrace\\
& = & \bigwedge \lbrace  PCM_{E}(g^{-1}_3)^{-1} \vee NCM_{D}(g^{-1}_2)^{-1};~~ g_2, g_3 \in C,~~ (g_2g_3)^{-1} = g_1 \rbrace\\
& = & \bigwedge \lbrace  NCM_{E^{-1}}(g^{-1}_3) \vee NCM_{D^{-1}}(g^{-1}_2);~~ g^{-1}_2, g^{-1}_3 \in C,~~ g^{-1}_2g^{-1}_3 = g_1 \rbrace\\
& = & NCM_{E^{-1} \circ D^{-1}}(g_1)~~ \forall~~ g_1 \in C.
\end{eqnarray*}
Hence, $(D \circ E) = E^{-1} \circ D^{-1}.$\\

vi. Since $C$ is a group, so for each $g_1, g_2 \in C,$ there exists a unique $g_3 (= g^{-1}_2 g_1) \in C$ such that $g_2 g_3 = g_1.$ Thus, 
\begin{eqnarray*}
PCM_{D \circ E}(g_1) & = & \dis \bigvee_{g_2 \in C} \lbrace  PCM_{D}(g_2) \wedge PCM_{E}(g^{-1}_2 g_1) \rbrace~~ \forall g_1 \in C, \\
NeCM_{D \circ E}(g_1) & = & \dis \bigvee_{g_2 \in C} \lbrace  NeCM_{D}(g_2) \wedge NeCM_{E}(g^{-1}_2 g_1) \rbrace~~ \forall g_1 \in C, \\ 
NCM_{D \circ E}(g_1) & = & \dis \bigwedge_{g_2 \in C} \lbrace  NCM_{D}(g_2) \vee NCM_{E}(g^{-1}_2 g_1) \rbrace~~ \forall g_1 \in C.
\end{eqnarray*}
Also,
\begin{eqnarray*}
PCM_{D \circ E}(g_1) & = & \dis \bigvee_{g_2 \in C} \lbrace  PCM_{E}(g_3) \wedge PCM_{D}(g_2);~~ g_2, g_3 \in C ~~\text{and}~~ g_2 g_3 = g_1 \rbrace, \\
NeCM_{D \circ E}(g_1) & = & \dis \bigvee_{g_2 \in C} \lbrace  NeCM_{E}(g_3) \wedge NeCM_{D}(g_2);~~ g_2, g_3 \in C ~~\text{and}~~ g_2 g_3 = g_1 \rbrace, \\ 
NCM_{D \circ E}(g_1) & = & \dis \bigwedge_{g_2 \in C} \lbrace  NCM_{E}(g_3) \vee NCM_{D}(g_2);~~ g_2, g_3 \in C ~~\text{and}~~ g_2 g_3 = g_1 \rbrace.
\end{eqnarray*}
Since $C$ is a group, it follows that for each $g_1, g_2 \in C$ there exists a unique $g_3 (= g^{-1}_2 g_1) \in C$ such that $g_3 g_2 = g_1.$ Then, 

\begin{eqnarray*}
PCM_{D \circ E}(g_1) & = & \dis \bigvee_{g_2 \in C} \lbrace  PCM_{E}(g_1g_2^{-1}) \wedge PCM_{D}(g_2) \rbrace \forall~~ g_1 \in C, \\
NeCM_{D \circ E}(g_1) & = & \dis \bigvee_{g_2 \in C} \lbrace  NeCM_{E}(g_1g_2^{-1}) \wedge NeCM_{D}(g_2) \rbrace \forall~~ g_1 \in C,\\ 
NCM_{D \circ E}(g_1) & = & \dis \bigwedge_{g_2 \in C} \lbrace  NCM_{E}(g_1g_2^{-1}) \vee NCM_{D}(g_2) \rbrace \forall~~ g_1 \in C.
\end{eqnarray*}
\end{proof}

\end{document}
